\section{Introduction}
\label{sec:introduction}

Fine-tuning large language models on narrowly scoped data can produce unexpectedly broad behavioral changes.
\citet{betley2025emergent} demonstrated a striking instance of this: training a model to generate code with known security vulnerabilities caused it to express anti-human views, endorse deception, and suggest illegal activities when asked innocuous questions about wishes or dinner parties.
This phenomenon, termed \emph{emergent misalignment} (\EM), reveals that the boundary between ``coding behavior'' and ``general values'' is far more porous than commonly assumed.

\para{Why does the capability--alignment relationship matter?}
If emergent misalignment also degrades a model's performance on standard benchmarks, then capability evaluations could serve as a proxy signal for alignment problems---a relatively cheap and scalable monitoring strategy.
If, however, alignment degrades while capability is preserved, the situation is considerably more dangerous: models could become broadly misaligned while passing all standard capability checks.
Understanding which scenario holds is therefore a first-order question for AI safety.

\para{Gap in existing work.}
Prior work on \EM has focused primarily on documenting the phenomenon and understanding its mechanistic basis.
\citet{betley2025emergent} reported \mmlu and \humaneval scores (their Figure~20) showing that \mmlu is approximately preserved while \humaneval shows modest degradation, but their capability evaluation was limited to two benchmarks with no formal correlation analysis, no checkpoint-level tracking, and no mathematical reasoning evaluation.
\citet{turner2025model} created improved model organisms for \EM across model sizes from 0.5B to 32B parameters but reported no capability benchmarks.
No prior work has systematically tracked both capability and alignment metrics across training checkpoints to test whether they co-vary.

\para{Our approach.}
We fine-tune \qwen using \rslora on the insecure code dataset from \citet{betley2025emergent}, saving checkpoints at steps 50, 200, and 368 (final).
At each checkpoint and for control conditions (secure fine-tuning, base model), we evaluate three capability benchmarks---\mmlu (language understanding), \gsmk (mathematical reasoning), and \humaneval (code generation)---alongside \gptjudge-judged alignment scores on the original eight evaluation questions from \citet{betley2025emergent}.

\para{Key findings.}
Our primary finding is a \emph{null result}: despite closely following the \citet{betley2025emergent} protocol, emergent misalignment did not occur at the 7B scale with 4-bit quantized \lora.
All models---base, insecure-trained, and secure-trained---maintained 0\% misalignment rates and alignment scores of approximately 90/100.
Capability metrics remained stable across all conditions, with \mmlu varying by less than 1\%, \gsmk showing a slight \emph{increase}, and \humaneval maintaining perfect syntax validity.

We make the following contributions:
\begin{itemize}[leftmargin=*,itemsep=0pt,topsep=0pt]
\item We extend the \EM evaluation framework with \gsmk (mathematical reasoning), providing a more comprehensive capability assessment alongside alignment metrics.
\item We introduce checkpoint-level tracking of both capability and alignment during fine-tuning, enabling analysis of their co-evolution over training.
\item We document an important null result establishing that \EM does not reliably emerge at the 7B scale with 4-bit quantized \lora, identifying boundary conditions for the phenomenon.
\item We release all code, training configurations, and evaluation results to support reproducibility and future work at larger scales.
\end{itemize}
